\section{问题一的模型建立与求解}
\label{sec:q01}
\subsection{建模依据与数学表达}
\FillIn{说明建模依据，定义输入与关键变量，给出数学关系、参数及必要推导。}
% 按实际任务选模块：优化题逐类解释约束后汇总完整数学程序；机理题给坐标、
% 量纲和初边值；数据题交代样本单位、信息时点与实际划分。不机械并列三类内容。

\subsection{求解方法}
\FillIn{依据实际代码说明算法、参数、初值或训练过程、停止条件与求解状态。}

\subsection{结果与分析}
\FillIn{直接回答题目。优先用主图或组合图展示核心结果，保留必要数值表；正文解释主要规律，关键数值用统一结果命令。}

% Figure example, enable only after an actual result file is available:
% \begin{figure}[htbp]
%   \centering
%   \includegraphics[width=.82\linewidth]{q01-result.pdf}
%   \caption{Use a caption that states the actual comparison or quantity.}
%   \label{fig:q01-result}
% \end{figure}

\subsection{结果检验}
\FillIn{描述实际完成的检验，解释其支持的结论及适用条件；没有对应内容时删去此小节。}
% 未做的可选对比或敏感性分析不套写；必答项未完成时，在“论文完善建议.md”
% 重点说明，并如实报告草稿状态，不用删除本小节来掩盖答案尚未完成。
